% Part: lambda-calculus
% Chapter: introduction
% Section: computable-lambda

\documentclass[../../../include/open-logic-section]{subfiles}

\begin{document}

\olfileid{lam}{int}{lrp}
\olsection{الدوال القابلة للحساب \usetoken{S}{lambda definable}}

\begin{thm}
\ollabel{thm:computable-lambda}
كل دالة جزئية قابلة للحساب !!{lambda definable}.
\end{thm}

\begin{proof}
نحتاج إلى إثبات أن كل دالة جزئية قابلة للحساب~$f$ !!{lambda defined}
بحد لامبدا~$F$. وبحسب مبرهنة كليني للصورة السوية، يكفي أن نثبت أن كل
دالة عودية بدائية !!{lambda defined} بحد لامبدا، ثم أن الدوال
!!{lambda definable} مغلقة تحت عمليات التركيب المناسبة والبحث غير
المحدود. ولإثبات أن كل دالة عودية بدائية !!{lambda defined} بحد لامبدا،
يكفي إثبات أن الدوال الابتدائية !!{lambda definable}، وأن الدوال الجزئية
!!{lambda definable} مغلقة تحت التركيب والعودية البدائية والبحث غير
المحدود.
\end{proof}

سنستعمل ترميزًا أكثر مألوفية لجعل بقية البرهان أيسر قراءة. فسنكتب،
مثلًا،~$M(x, y, z)$ بدلًا من~$Mxyz$. ومع أن هذا موحٍ، ينبغي تذكر أن
الحدود في حساب لامبدا غير المنمّط ليست لها مواضع ثابتة؛ ولذلك يصح، للحد
نفسه~$M$، أن نكتب~$M(x,y)$ أو~$M(x,y,z,w)$. غير أن هذا الترميز يشير
إلى أننا نتعامل مع~$M$ بوصفه دالة في ثلاثة متغيرات، ويوضح مقاصد
التعريفات. وبالمثل سنقول «عرّف~$M$ بـ$M(x,y,z) = \dots$» بدلًا من
«عرّف~$M$ بـ$M = \lambd[x][\lambd[y][\lambd[z][\dots]]]$».

\end{document}
